\documentclass[12pt,a4paper]{article}

\usepackage[utf8]{inputenc}
\usepackage[T1]{fontenc}
\usepackage{amsmath,amssymb,amsthm}
\usepackage{booktabs}
\usepackage{array}
\usepackage{graphicx}
\usepackage{natbib}
\usepackage[margin=1in]{geometry}
\usepackage[hyphens]{url}
\usepackage{hyperref}
\usepackage{setspace}
\usepackage{lineno}
\Urlmuskip=0mu plus 1mu\relax
\onehalfspacing
\emergencystretch=2em

\newtheorem{assumption}{Assumption}
\newtheorem{theorem}{Theorem}
\newtheorem{proposition}{Proposition}
\newtheorem{corollary}{Corollary}
\newtheorem{remark}{Remark}
\newcommand{\E}{\mathbb{E}}
\newcommand{\Var}{\mathbb{V}\mathrm{ar}}
\newcommand{\Cov}{\mathbb{C}\mathrm{ov}}
\newcommand{\Corr}{\mathbb{C}\mathrm{orr}}
\newcommand{\SR}{\mathrm{SR}}
\newcommand{\UVIF}{\mathrm{UVIF}}
\newcommand{\Neff}{N_{\mathrm{eff}}}
\DeclareMathOperator{\sign}{sign}

\title{Sharpe-ratio variance inflation under cross-sectional and serial dependence in trading panels}
\author{Dhananjay S. Chauhan\\
Independent Researcher, India\\
\texttt{dschauhan08.me@gmail.com}\\
ORCID: \url{https://orcid.org/0009-0003-1049-2213}}
\date{}

\begin{document}
\linenumbers
\maketitle

\noindent\textbf{Correspondence:} Dhananjay S. Chauhan, Independent
Researcher, India.\\
\textbf{Email:} \texttt{dschauhan08.me@gmail.com}\\
\textbf{ORCID:} \url{https://orcid.org/0009-0003-1049-2213}.

\begin{abstract}
Trading-strategy evidence is often produced from symbol-date panels,
although the economically tradable object is a date-level portfolio
return. Treating selected rows as independent can therefore overstate
evidence when same-date rows share market, factor, liquidity, or
volatility shocks and when the resulting date-level returns are serially
dependent. This paper develops a dependence-aware inference framework
for trading-panel evidence. The main theoretical object is the Sharpe
Unified Variance Inflation Factor, defined as the ratio of the long-run
delta-method variance of the date-level Sharpe estimator to the
row-pooled IID delta-method variance. The framework combines date
aggregation, HAC-delta inference, dependent bootstrap resampling,
researcher-menu correction, same-date placebo tests, factor-alpha
diagnostics, and turnover-scaled cost checks. Public benchmark and
stress-test evidence from Kenneth French and AQR data shows that
canonical momentum survives the dependence-aware boundary, while weaker
panel claims can lose much of their nominal row-level information after
same-date redundancy and date-level inference are imposed. Monte Carlo
experiments show that row-naive inference can substantially overreject
under dependent nulls, while date-level HAC, block-bootstrap,
stationary-bootstrap, and Romano-Wolf procedures remain closer to
nominal size. The results provide a reproducible claim boundary for
trading-panel evidence and show why statistical significance produced
by multiplying dependent rows should not be interpreted as evidence of
a tradable edge.
\end{abstract}

\noindent\textbf{Keywords:} Sharpe ratio; financial econometrics; panel data; cross-sectional
dependence; serial dependence; HAC inference; dependent bootstrap; multiple
testing; trading strategy evaluation; information redundancy

\noindent\textbf{JEL Codes:} C12; G11; G17

\section{Introduction}

\subsection{Problem: The wrong sampling boundary}

Prediction-based trading studies commonly begin with a panel
\(\mathcal P=\{(t,i,s_{t,i},r_{t+1,i})\}\), where \(t\) indexes dates,
\(i\) indexes assets or portfolios, \(s_{t,i}\) is a signal, and
\(r_{t+1,i}\) is the next-period return.  One evaluation failure mode is to
evaluate a positive edge using all selected symbol-date rows as though
they were independent.  This turns a panel with \(T\) dates and \(n_t\)
selected rows per date into a nominal sample of \(\sum_t n_t\), even
when same-date rows share common shocks.  The paper does not claim that
this is best practice among careful empirical asset-pricing researchers;
it uses the row-level calculation as a baseline reference for bias
quantification and as a concrete failure mode for automated strategy
searches.

The problem is not only a matter of notation.  A row-pooled \(t\)-test
implicitly treats the cross-section on a date as many independent draws
from the same experiment.  A trading strategy, however, is marked to
market by date.  If a strategy selects several portfolios on the same
date, the economically relevant object is the date-level portfolio
return after aggregation.  Cross-sectional dependence determines how
much extra information the additional selected rows actually contain,
and serial dependence determines the long-run variance of the date
series.
Standard clustering techniques, including the finance-panel guidance of
\citet{petersen2009estimating}, address standard errors for linear
regression coefficients or means.  They do not by themselves define a
complete inference record for nonlinear performance metrics such as the
Sharpe ratio when the same data contain both cross-sectional shocks and
serial memory.  The missing object is a single variance-inflation
accounting that starts from the row-pooled benchmark, moves through the
date-return boundary, and ends at the long-run variance of the Sharpe
functional.

Two terms are used throughout the paper.  The Sharpe Unified Variance
Inflation Factor, or Sharpe UVIF, is
\[
  \UVIF_{\SR}
  =
  \frac{\Var_{LR}(\widehat{\SR}_{date})}
       {\Var_{IID,row}(\widehat{\SR}_{row})}.
\]
It is a design-effect analogue for the nonlinear Sharpe ratio: values
above one mean that row-pooled IID inference understates Sharpe
uncertainty, extending the clustered-sampling logic of
\citet{kish1965survey} from means to the Sharpe functional.  The
Effective Information Rate (EIR) is a diagnostic
fraction of row-level information that remains nonredundant after
same-date dependence is recognized.  Under the equicorrelation
approximation, \(\mathrm{EIR}=1/\{1+(n_t-1)\rho\}\); for example,
\(\mathrm{EIR}=0.11\) means roughly 89 percent of nominal row-level
information is redundant.

\subsection{Solution: Dependence-aware inference}

The paper's estimator-level claims are deliberately conservative.  HAC,
bootstrap, Romano-Wolf, and portfolio-level testing are standard tools.
The contribution is the way they are assembled at the correct sampling
boundary and the Sharpe-specific variance accounting that makes the cost
of row pooling explicit.  First, the paper gives a compact
trading-panel version of the familiar Moulton-like design-effect calculation
for the standard error of a row mean.  Second, it derives a Unified
Variance Inflation Factor for the Sharpe ratio, replacing a mean-only
inflation factor with the delta-method variance of the nonlinear
functional \(f(\mu,m_2)\).  Third, it turns standard econometric tools
into a pre-specified two-channel decision record: date aggregation
addresses cross-sectional clustering, and HAC-delta plus moving-block or
stationary bootstrap inference addresses serial dependence in the
date-return series.  Fourth, it reports positive-edge claims alongside
data-snooping, factor-alpha, placebo, stationarity, and cost diagnostics.

The empirical evidence is framed as a public benchmark and stress-test
campaign, not as a toy positive-control example.  It asks
which familiar momentum, value, quality, beta, and size/book-to-market
claims survive once the sampling boundary, dependence, bootstrap
robustness, researcher menu, placebo, factor-alpha, and cost boundaries
are imposed.  The public Kenneth French and AQR data are rendered into
the manuscript from
production CSV artifacts.
The conservative empirical question is whether apparently significant
trading claims are artifacts of row-level sampling and disappear after
date aggregation, dependence adjustment, and researcher-menu correction.
The same code also runs a Monte Carlo study comparing row-naive testing
with date-level bootstrap, HAC-delta, block bootstrap, stationary
bootstrap, and joint resampling decisions.

\section{Related work}

This paper sits at the intersection of three literatures that are often
discussed separately.  Clustered-panel econometrics explains why row
counts overstate information, Sharpe-ratio and dependent-bootstrap
methods define the date-level sampling law, and data-snooping work in
the factor zoo explains why a single winning parameter choice is not a
research claim.  The protocol combines those literatures into one
pre-specified decision record.

The sampling-boundary argument follows survey and grouped-error design
effects \citep{kish1965survey,moulton1986random} and finance-panel
warnings about clustered standard errors \citep{petersen2009estimating}.
Sharpe-ratio inference follows the time-series and HAC literature
\citep{lo2002statistics,ledoit2008robust,opdyke2007comparing,andrews1991heteroskedasticity,kiefer2005new};
dependent-bootstrap and researcher-menu corrections follow block
bootstrap, Reality Check, and Romano-Wolf procedures
\citep{kunsch1989jackknife,politis1994stationary,lahiri2003resampling,white2000reality,romano2005stepwise}.
The empirical interpretation is tied to the factor-zoo, false-discovery,
factor-alpha, anomaly-decay, and trading-cost literatures
\citep{harvey2016and,harvey2020false,harvey2020census,bailey2014deflated,holm1979simple,fama1993common,carhart1997persistence,fama2010luck,mclean2016does,almgren2001optimal,hasbrouck2009trading,frazzini2015trading}.

\section{Materials and methods}
\label{sec:method}

The protocol is a dependence-aware inference framework rather than a
single test.  Its first channel is cross-sectional: all evidence is
translated from selected symbol-date rows into one economically
tradeable date return, so same-date clusters cannot create sample size by
row multiplication.  Its second channel is serial: all
Sharpe decisions are made on the date-return series using HAC-delta and
dependent-bootstrap procedures.  A candidate must survive both channels
before the paper asks the separate researcher-menu, placebo, factor-alpha,
and cost questions.

Operationally, the audit has three reported scales.  Row-naive statistics
are retained only as the baseline to be audited.  Date-level Sharpe
statistics are the objects used for confirmatory inference.  Sharpe UVIF
and EIR translate the gap between those two scales into variance and
information terms, respectively.

\subsection{Signals and executability}

For lookback \(L\) and skip \(q\), the skipped compounded signal is
\begin{equation}
  s_{t,i}^{(L,q)} =
  \prod_{j=t-L-q+1}^{t-q} (1+r_{j,i}) - 1 .
  \label{eq:signal}
\end{equation}
The bounds in \eqref{eq:signal} are ascending.  With \(q\ge1\), the
signal excludes the outcome return and is executable at date \(t\) when
all signals in the available universe \(A_t\) are known before the
trading decision.  This simultaneity condition is automatic for the
public Kenneth French portfolio panel because all signals are computed
from lagged public returns.  It is a required condition for applying the
protocol to any other prediction panel.

The sign convention is \(\sign(0)=0\).  Zero-signal and missing-signal
rows receive zero economic weight and are excluded from threshold
selection.  The signed one-period row return is
\begin{equation}
  y_{t+1,i}^{(L,q)} = \sign(s_{t,i}^{(L,q)})r_{t+1,i}.
\end{equation}

The confidence score is the within-date percentile rank of
\(|s_{t,i}^{(L,q)}|\).  This percentile score is a baseline screening
heuristic, not an optimal weighting theorem.  It is scale-free,
transparent, and invariant to monotone transformations of the signal
magnitude.  Z-scores, volatility-scaled scores, and learned weights are
valid alternatives only if they are declared before the confirmatory run
or included in the researcher menu.  The protocol's main object is the
sampling boundary, not weighting optimality.  With \(N_t\) finite
non-zero signals and average
midrank \(R_{t,i}\),
\begin{equation}
  c_{t,i}^{(L,q)}=
  \begin{cases}
    \dfrac{R_{t,i}-1}{N_t-1}, & N_t>1,\\
    1, & N_t=1 .
  \end{cases}
  \label{eq:confidence}
\end{equation}
For dates with \(N_t=1\), the convention \(c_{t,i}=1\) represents a
single-asset conviction rather than an undefined percentile rank.  Dates
with \(N_t=0\) are inactive for that menu element.  Ties use average
midranks.  When \(N_t=2\), the available confidence scores are
\((0,1)\), or \((0.5,0.5)\) under a tie, so high threshold decisions are
coarse.  The empirical tables therefore report selected counts by
threshold rather than treating the threshold scale as equally granular
across dates.

\subsection{Date returns and effective sample size}

For threshold \(\tau\), selected set
\(\mathcal S_t(L,q,\tau)=\{i:c_{t,i}^{(L,q)}\ge\tau\}\), and optional
non-negative weights \(a_{t,i}\), the date return is
\begin{equation}
  \bar r_{t+1}(L,q,\tau) =
  \frac{\sum_{i\in\mathcal S_t} a_{t,i}y_{t+1,i}^{(L,q)}}
       {\sum_{i\in\mathcal S_t} a_{t,i}} .
  \label{eq:date-return}
\end{equation}
The denominator is the total active economic weight on date \(t\), not a
generic row count.  Under the primary equal-weighted evaluation
\(a_{t,i}=1\), it reduces to the number of selected rows.  Confidence-weighted
and signal-magnitude-weighted returns are sensitivity checks.

For descriptive diagnostics, the artifacts report a positive-sequence
effective sample size for the finite date-return series.  The
autocorrelation sum is truncated at the first negative sample
autocorrelation, up to a fixed maximum lag.  This is not a substitute for
HAC or bootstrap inference; it is a compact diagnostic for the loss of
information from serial dependence and is not used as a decision rule.

\subsection{Sharpe inference}

The protocol separates estimation from testing.  It reports two-sided
95 percent intervals for \(\widehat{\SR}\) and one-sided positive-edge
\(p\)-values for \(H_0:\SR\le0\).  Bootstrap inference resamples the
date-return series by IID date, moving block, or stationary bootstrap.
The positive-edge bootstrap \(p\)-value uses one-draw smoothing:
\begin{equation}
  p^+ = \frac{1+\#\{\widehat{\SR}_b^*\le 0\}}{B+1}.
\end{equation}

For HAC inference, reindex the finite realized date-return observations
as \(R_u\), \(u=1,\ldots,T_R\), where \(R_u\) is the strategy excess
return for the decision date after subtracting the applicable risk-free
rate when the source is not already a long-short excess-return series.
Let \(g_u=(R_u,R_u^2)'\),
\(\hat\mu=T_R^{-1}\sum_u R_u\),
\(\hat m_2=T_R^{-1}\sum_u R_u^2\), and
\(\hat v=\hat m_2-\hat\mu^2\).  The denominator uses the \(T_R\)-moment
variance, not the \(T-1\) bias correction, because it is the plug-in
moment entering the smooth functional
\[
  f(\mu,m_2)=\sqrt{252}\frac{\mu}{(m_2-\mu^2)^{1/2}}.
\]
Following \citet{ledoit2008robust}, the Sharpe ratio is treated as a
smooth functional \(f(\theta)\) of the moment vector
\(\theta=\E[g_t]\).  On sets where \(m_2-\mu^2\) is bounded away from
zero, this map is Hadamard differentiable.  The asymptotic normality of
\(\sqrt{T_R}(\widehat{\SR}-\SR)\) follows from the functional delta method
provided the long-run covariance matrix \(\Omega\) is consistently
estimated by a kernel-based HAC estimator.
Evaluated at \((\hat\mu,\hat m_2)\),
\[
  \widehat{\SR}=\sqrt{252}\frac{\hat\mu}{\sqrt{\hat v}},
  \qquad
  \nabla f(\hat\mu,\hat m_2) =
  \left(
  \sqrt{252}\frac{\hat m_2}{\hat v^{3/2}},
  -\frac{1}{2}\sqrt{252}\frac{\hat\mu}{\hat v^{3/2}}
  \right)' .
\]
With \(\widehat{\Omega}\) the Bartlett-kernel HAC estimate of the
long-run covariance of \(g_u\), the HAC-delta variance is
\begin{equation}
  \widehat{\Var}(\widehat{\SR})
  =\nabla f(\hat\mu,\hat m_2)'\widehat{\Omega}
   \nabla f(\hat\mu,\hat m_2)/T_R .
  \label{eq:hac-delta}
\end{equation}
Explicitly, for centered moments \(\tilde g_u=g_u-\bar g\),
\[
  \widehat{\Gamma}_k
  = T_R^{-1}\sum_{u=k+1}^{T_R}\tilde g_u\tilde g_{u-k}',
  \qquad
  \widehat{\Omega}
  = \widehat{\Gamma}_0
    +\sum_{k=1}^{K_T} \left(1-\frac{k}{K_T+1}\right)
      \left(\widehat{\Gamma}_k+\widehat{\Gamma}_k'\right).
\]
This is a covariance matrix for the joint process
\((R_u,R_u^2)\), so the off-diagonal terms carry the
long-run covariance between the sample mean and second moment into the
Sharpe-ratio delta method.
The bandwidth \(K_T\) is a tuning choice and may depend on \(T\).  The
implementation reports the
selected bandwidth and uses a documented automatic Bartlett bandwidth as
the default; important empirical claims should be checked against
reasonable nearby bandwidths and against the block-bootstrap intervals
rather than relying on a single HAC lag choice.
The technical appendix reports prewhitened HAC and fixed-fraction HAC
sensitivity checks.

Because financial date returns can be heavy-tailed, the HAC-delta
calculation is not allowed to stand alone.  The manuscript reports
bootstrap \(p\)-values and intervals adjacent to HAC estimates; when the
bootstrap evidence is materially weaker than the HAC evidence, the
candidate is treated as robustness-fragile rather than as a clean
positive-edge finding.

\subsection{Researcher menus, placebos, and factor alpha}

For a fixed menu \(\mathcal M\) of lookbacks and thresholds, the primary
joint adjustment is Romano-Wolf stepdown max-\(t\) resampling over
\(\{R_u(m):m\in\mathcal M\}\).  Holm-adjusted marginal bootstrap
\(p\)-values are secondary.  The protocol also reports White's Reality Check
for the maximum mean return and the Deflated Sharpe Ratio diagnostic of
\citet{bailey2014probability}.
In the public campaign, the menu is the finite registry-defined set
\[
  \mathcal M=\{(c,\theta_c,\tau):c\in\mathcal C,\ \tau\in\mathcal T_c\},
\]
where \(\theta_c\) denotes the candidate's frozen signal definition.
For panel candidates, \(\mathcal T_c=\{0.0,0.3,0.5,0.7,0.9\}\) and
any lookback and skip values are fixed in \path{code/campaign_registry.json}
rather than selected after seeing results.  Pre-aggregated AQR factor-series
candidates use the singleton threshold menu \(\mathcal T_c=\{0.0\}\);
because they do not contain row-level constituent signals, same-date
panel permutation is not applicable to them.
The reason is mechanical.  If \(N\) independent strategies with
zero true Sharpe are searched over a sample of length \(T\), the maximum
observed daily Sharpe is on the order of \(\sqrt{2\log(N)/T}\) before any
economic signal exists.  Real researcher menus are usually correlated
rather than independent, so the protocol uses Romano-Wolf stepdown to
control familywise error while preserving the joint dependence among
menu elements.

The fixed-menu condition is a confirmatory convention, not a claim that
research is never iterative.  In practice, exploratory work should be
separated from the final evaluation by freezing the signal definitions,
lookbacks, thresholds, filters, cost assumptions, and benchmark model
before the confirmatory run.  If discarded variants influenced the final
choice, they belong in the researcher menu or the exercise should be
reported as exploratory.  A stronger design uses a locked holdout period
or an independently refreshed sample after the menu is frozen.

The public placebo is a same-date permutation design.  For each date,
the vector of signals is randomly permuted across portfolios, preserving
the date's cross-sectional signal distribution and missingness while
breaking the signal-return pairing.  The signed next-day returns,
confidence ranks, selected counts, turnover, and date aggregation are
then recomputed.  Repeating this procedure gives a permutation null for
the observed public momentum panel.

This placebo is a test under a within-date exchangeability hypothesis:
under the null, the assignment of signals \(s_{t,i}\) to next-period
returns \(r_{t+1,i}\) is exchangeable within each date \(t\).  If signals
are structurally tied to size, sector, liquidity, country, exchange,
beta, or other common-factor state variables, an unrestricted same-date
permutation can break the wrong null.  In those cases the protocol should
permute within pre-specified blocks, residualize returns against the
structural variables before permutation, or mark the placebo as not
applicable.  The permutation gate is therefore a design-specific
diagnostic rather than a universal validity guarantee.

The factor-alpha test estimates
\begin{equation}
  R_t = \alpha + \beta_M(MKT_t-RF_t)+\beta_S SMB_t
  +\beta_H HML_t+\beta_U MOM_t+\varepsilon_t ,
  \label{eq:alpha}
\end{equation}
using HAC standard errors.  The null is \(H_0:\alpha\le0\).  Including
momentum is essential because the public signal is itself momentum-like.

\subsection{Turnover-scaled costs}

For the equal-weighted threshold portfolio, define target weights
\begin{equation}
  w_{t,i} =
  \frac{\sign(s_{t,i})\mathbf 1\{i\in\mathcal S_t\}}
       {|\mathcal S_t|},
  \qquad w_{t,i}=0 \text{ if } |\mathcal S_t|=0 .
\end{equation}
Daily turnover is the half-\(L^1\) distance between adjacent target
weight vectors,
\begin{equation}
  TO_t = \frac{1}{2}\sum_i |w_{t,i}-w_{t-1,i}|,
  \qquad
  \widehat{TO}=T_{TO}^{-1}\sum_t TO_t ,
  \label{eq:turnover}
\end{equation}
where the average is over adjacent dates with at least one active side.
For cost \(c\) per full rebalance in decimal units, the cost-adjusted
Sharpe ratio is
\begin{equation}
  \widehat{\SR}_{net}(c)
  = \frac{\hat\mu-\widehat{TO}\,c}{\hat\sigma}\sqrt{252}.
  \label{eq:net-sr}
\end{equation}
This is a first-order cost stress.  It assumes linear cost per unit of
turnover and holds the return volatility fixed after subtracting costs.
Those assumptions are reasonable for a screening evaluation of diversified,
low-frequency panels, and the resulting cost drag should be read as a
lower bound when liquidity demand is material.  They are not a full execution model for
high-frequency, capacity-constrained, or illiquid strategies.  Such
applications require spread, impact, latency, borrow, and capacity
models before any tradability claim.  The protocol establishes a
statistical lower bound for viability; it does not account for
liquidity-driven alpha decay at larger assets under management.
The empirical artifacts also report the break-even cost in basis points,
defined by \(\hat\mu-\widehat{TO}c=0\).  A square-root impact capacity
extension in the spirit of \citet{almgren2001optimal} and the broader
market-impact literature \citep{bouchaud2008markets} is provided in the
technical appendix.  The public sorted-portfolio examples do not contain
average daily volume or order-size information, so they do not support a
capacity claim.

\subsection{Composite viability decision}

The protocol reports component diagnostics, but the final claim boundary is
deliberately stricter than any single statistic.  The gates are a
lexicographic filter: a candidate must first clear the date-level
statistical boundary, then the dependent-bootstrap robustness boundary,
then the researcher-menu and placebo integrity boundaries, and finally
the economic cost boundary.  A candidate is called viable only if it
clears a pre-specified composite gate:
\[
  p_{\mathrm{boot}}\le \alpha,\qquad
  p_{\mathrm{perm}}\le \alpha,\qquad
  p_{\mathrm{HAC}}\le \alpha,\qquad
  p_{\mathrm{RW}}\le \alpha,\qquad
  \widehat{\SR}_{net}(c)>0 .
\]
This is a conservative policy rule, not an optimal weighted test or a
claim that the gates exhaust all possible failure modes.  The gates
are conjunctive because each one addresses a different reason a trading
claim can fail: bootstrap-fragile Sharpe evidence, a placebo-consistent
signal, dependence-adjusted insignificance, researcher-menu selection,
or nonpositive net economics.
Applications with different decision costs may pre-specify different
gates, but they should not choose the gate sequence after seeing the
results.
The reported artifacts make the gate threshold reproducible: in addition to
the conventional \(\alpha=0.05\), a reader can recompute pass counts at
stricter levels such as \(\alpha=0.01\) without changing the statistical
objects being tested.
The same-date permutation gate is a panel placebo: it is applicable only
when row-level signals can be permuted within date.  Pre-aggregated
factor series can still be economically interesting and can pass
HAC-delta or Romano-Wolf tests, but they do not contain the row-level
signal panel needed for this placebo.  Such cases are reported as not
applicable for the permutation gate rather than as passed panel evaluations.

The stages have different interpretations.  The permutation gate asks
whether the signal-return pairing beats a same-date placebo.  The HAC
gate asks whether the resulting date-return Sharpe is positive after
serial dependence and heteroskedasticity are accounted for.  The
bootstrap gate asks whether that conclusion survives resampling of the
date-return series.  The Romano-Wolf gate asks whether the claim survives
the fixed researcher menu.  The cost gate asks whether the edge remains
economically usable after turnover-scaled frictions.  Factor-alpha,
holdout, subperiod, and stationarity diagnostics are reported alongside
these gates to prevent a single favorable statistic from becoming the
claim.

\subsection{Sampling-boundary theory}
\label{sec:theory}

\begin{assumption}[Equicorrelated row model]
For each date \(t\), selected row returns satisfy
\[
  y_{t,i}=\mu+\lambda_t+\epsilon_{t,i},
\]
where \(\lambda_t\) is a common date shock, \(\epsilon_{t,i}\) is
idiosyncratic, \(\lambda_t\) is uncorrelated with \(\epsilon_{t,i}\),
\(\Var(\lambda_t)=\sigma_\lambda^2\), and
\(\Var(\epsilon_{t,i})=\sigma_\epsilon^2\).  The selected count is
\(n_t=\bar n\) for the displayed derivation, and dates are independent.
\end{assumption}

Under this model,
\[
  \sigma^2=\Var(y_{t,i})=\sigma_\lambda^2+\sigma_\epsilon^2,
  \qquad
  \rho=\Corr(y_{t,i},y_{t,j})
       =\frac{\sigma_\lambda^2}{\sigma_\lambda^2+\sigma_\epsilon^2}
       \quad (i\ne j).
\]
Thus \(\rho\) is the intraclass correlation induced by the common date
shock; economically, it is the fraction of row-return variance explained
by shocks shared by same-date selections.
In empirical diagnostics, \(\rho\) is estimated after the candidate
threshold is fixed.  For each panel candidate, the selected signed row
returns are aligned by date and asset or portfolio, and
\(\hat\rho\) is the average off-diagonal sample correlation among active
row-return series with overlapping dates; \(\bar n\) is the average
selected count across active dates.  When the selected set varies over
time, this equicorrelation estimate is a descriptive compression of the
same-date covariance structure, not an identifying assumption for the
confirmatory HAC or bootstrap tests.

\begin{proposition}[Mean-return design effect]
Under the equicorrelated row model, the ratio of the naive row-pooled
reported standard error of the mean to the date-level standard error is
\[
  \frac{se_{row,naive}}{se_{date,true}}
  =
  \left\{1+(\bar n-1)\rho\right\}^{-1/2}.
\]
For \(\rho>0\) and \(\bar n>1\), the naive row-level report is too small.
\end{proposition}

This calculation is not new as a sampling result; it is the
Moulton-like special case of the paper's UVIF logic, written for signed
trading-panel rows and date-level portfolio returns.  To avoid conflict
with the researcher menu notation \(\mathcal M\), write the
cross-sectional variance design effect as
\[
  \UVIF_{\mu,xs}(\bar n,\rho)=1+(\bar n-1)\rho.
\]
The square root is reported when the object is a standard-error or
\(t\)-statistic adjustment:
\[
  \UVIF_{\mu,xs}^{1/2}(\bar n,\rho)=
  \sqrt{1+(\bar n-1)\rho}.
\]
Thus \(\UVIF\) denotes a variance ratio, while \(\UVIF^{1/2}\) denotes
the associated standard-error multiplier.  This mean-return calculation
is expository: it makes the scale of understatement visible before the
paper turns to HAC and bootstrap decisions on the realised date-return
series.
The magnitude can be economically large.  The relevant inflation factor
for the true standard error relative to the naive row report is
\(\UVIF_{\mu,xs}^{1/2}(\bar n,\rho)\).  With \(\bar n=50\), increasing same-date
correlation from \(\rho=0.10\) to \(\rho=0.60\) raises this factor from
2.43 to 5.51.  Thus the largest damage occurs precisely in regimes where
many names are selected and common shocks dominate idiosyncratic noise,
which is the pattern one expects during market stress.

\noindent The proof is given in the technical appendix.

\begin{remark}[Heterogeneous factor covariance]
The equicorrelated model is a closed-form illustration, not a restriction
imposed by the protocol.  Let a same-date signed-return vector satisfy
\[
  y_t=\mu\mathbf 1 + B f_t + u_t,
\]
where \(B\) contains asset-specific loadings, \(\Var(f_t)=\Sigma_f\),
and \(\Var(u_t)=D_t\).  Then the same-date covariance matrix is
\[
  \Sigma_t = B\Sigma_f B' + D_t,
\]
with heterogeneous pairwise correlations whenever loadings or
idiosyncratic variances differ.  For any displayed portfolio weights
\(w_t\), the relevant one-date variance is \(w_t'\Sigma_t w_t\), not the
row-independent variance.  The protocol avoids estimating the full
\(\Sigma_t\) as a prerequisite for inference by moving the evidence to
the realised date-return series and then using HAC or dependent
bootstrap inference on that series.  The design-effect calculation is therefore best read
as the transparent special case that explains why row pooling fails.
\end{remark}

\subsection{Unified variance inflation for the Sharpe ratio}

The mean-return design effect is useful but incomplete because the
reported performance metric is usually the Sharpe ratio.  The Sharpe
ratio is a nonlinear functional of the first two moments, so its
variance inflation depends on the covariance of \(R_t\) and \(R_t^2\),
not only on the variance of \(R_t\).  This motivates the Unified
Variance Inflation Factor:
\begin{equation}
  \UVIF_{\SR}
  =
  \frac{\Var_{LR}(\widehat{\SR}_{portfolio})}
       {\Var_{naive}(\widehat{\SR}_{row})}.
  \label{eq:uvif-definition}
\end{equation}
The numerator is the long-run delta-method variance at the
date-return boundary.  The denominator is the row-pooled delta-method
variance that would be reported if the selected rows were treated as
independent observations.
Equation \eqref{eq:uvif-definition} is therefore the primary definition:
Sharpe UVIF is a variance inflation factor.  Its square root is the
standard-error multiplier for Sharpe inference.  Under exchangeability
of the linearized Sharpe influence function, the cross-sectional part of
this variance ratio reduces to the same \(1+(\bar n-1)\rho\) design
effect shown above, with \(\rho\) applied to the linearized Sharpe
influence rather than to raw returns.

\begin{proposition}[Sharpe UVIF decomposition]
Let \(R_t\) be the date-level portfolio excess return and
\(G_t=(R_t,R_t^2)'\).  Define
\(\theta_R=(\mu_R,m_{2,R})'=\E[G_t]\),
\(v_R=m_{2,R}-\mu_R^2>0\), and
\[
  d_R=\nabla f(\theta_R)
  =
  \left(
  \sqrt{252}\frac{m_{2,R}}{v_R^{3/2}},
  -\frac{1}{2}\sqrt{252}\frac{\mu_R}{v_R^{3/2}}
  \right)' .
\]
Let \(\Gamma_k^R=\Cov(G_t,G_{t-k})\) and
\(\Omega_R=\Gamma_0^R+\sum_{k\ge1}(\Gamma_k^R+\Gamma_k^{R\prime})\).
For the selected row return \(Y_{t,i}\), let
\(H_{t,i}=(Y_{t,i},Y_{t,i}^2)'\),
\(\theta_Y=\E[H_{t,i}]\), \(d_Y=\nabla f(\theta_Y)\), and
\(\Psi_Y=\Var(H_{t,i})\).  With a constant selected count
\(\bar n\), the first-order Sharpe variance inflation relative to a
row-pooled IID report is
\begin{equation}
  \UVIF_{\SR}
  =
  \underbrace{
  \frac{\bar n\, d_R'\Gamma_0^R d_R}
       {d_Y'\Psi_Y d_Y}
  }_{\text{cross-sectional UVIF component}}
  \times
  \underbrace{
  \frac{d_R'\Omega_R d_R}
       {d_R'\Gamma_0^R d_R}
  }_{\text{Sharpe long-run serial factor}} .
  \label{eq:uvif-decomposition}
\end{equation}
Moreover, the serial factor can be written as
\[
  1+2\sum_{k=1}^{\infty}\psi_k^{\SR},
  \qquad
  \psi_k^{\SR}
  =
  \frac{d_R'\Gamma_k^R d_R}{d_R'\Gamma_0^R d_R},
\]
and the scalar long-run variance entering the numerator is
\begin{equation}
  d_R'\Omega_R d_R
  =
  \frac{252}{v_R^3}
  \left[
    m_{2,R}^2\Omega_{11}
    -\frac{\mu_R m_{2,R}}{2}(\Omega_{12}+\Omega_{21})
    +\frac{\mu_R^2}{4}\Omega_{22}
  \right],
  \label{eq:uvif-fourth-moment}
\end{equation}
where \(\Omega_{ab}\) denotes the \((a,b)\) element of \(\Omega_R\).
\end{proposition}

\noindent The proof is given in the technical appendix.

\begin{remark}[Connection to Moulton-like design effects]
If the linearized row Sharpe influence
\(\zeta_{t,i}^Y=d_Y'(H_{t,i}-\theta_Y)\) is exchangeable within date
with intraclass correlation \(\rho_\zeta\), and the date-level
linearization is the equal-weighted average of those row influences, the
cross-sectional factor in \eqref{eq:uvif-decomposition} reduces to
\(1+(\bar n-1)\rho_\zeta\).  The mean-return design effect is the
special case in which the influence variable is \(Y_{t,i}-\mu\).  The
Sharpe case replaces raw return correlation by correlation of the
linearized nonlinear-performance influence, which is why skewness,
kurtosis, and volatility clustering enter the variance inflation.
\end{remark}

\subsection{Higher-moment UVIF and non-normality}

The two-moment Sharpe UVIF already depends on fourth moments through the
variance of \(R_t^2\).  If a venue or application uses an adjusted Sharpe
ratio that explicitly corrects for skewness and kurtosis, the same
delta-method construction extends the moment vector from
\((R_t,R_t^2)'\) to \((R_t,R_t^2,R_t^3,R_t^4)'\).  The resulting
higher-moment UVIF is the ratio of the adjusted-Sharpe long-run variance
at the date-return boundary to the corresponding row-pooled IID variance.
It is reported as a tail-risk diagnostic, not as an additional rejection
gate.  The full expression is in the technical appendix.

\subsection{Effective information rate}

The UVIF can also be read as an information-redundancy statement.  Using
standard mutual-information notation \citep{cover2006elements}, define
the Effective Information Rate (EIR) for a date with \(n_t\) active rows
as
\begin{equation}
  \mathrm{EIR}_t
  =
  \frac{I(S_t;R_t)}
       {n_t\,\bar I_t},
  \qquad
  \bar I_t=n_t^{-1}\sum_{i=1}^{n_t}I(s_{t,i};r_{t+1,i}),
  \label{eq:eir-definition}
\end{equation}
where \(S_t\) denotes the vector of same-date signals and \(R_t\) is the
date-level portfolio return.  Under the same equicorrelated Gaussian
linear approximation that gives the mean-return UVIF, the information
rate relative to \(n_t\) independent rows is
\begin{equation}
  \mathrm{EIR}_t(\rho)
  \approx
  \frac{1}{1+(n_t-1)\rho},
  \qquad \rho\ge0.
  \label{eq:eir-equicorr}
\end{equation}
As \(\rho\to1\), each additional same-date row contributes mostly
redundant information and the rate per row falls toward \(1/n_t\).  This
is the information-theoretic version of the audit's core claim:
row-level pooling can report sample size that is not present in the
economic experiment.  In empirical tables, EIR is reported only as a
diagnostic approximation and is not used as a rejection rule.

\begin{assumption}[Fixed-menu bootstrap validity]
For each fixed menu element \(m\), the date-return series
\(\{\bar r_t(m)\}\) is strictly stationary and alpha-mixing with moments
of order \(4+\delta\).  The long-run covariance of
\((\bar r_t,\bar r_t^2)\) is finite and non-degenerate, and
\(\Var(\bar r_t)\ge\eta>0\).  The block length satisfies
\(\ell\to\infty\) and \(\ell/T\to0\).  The researcher menu
\(\mathcal M\) is fixed and finite.
\end{assumption}

\begin{corollary}[Fixed-menu validity statement]
Under the fixed-menu bootstrap-validity assumption, the HAC-delta and
block-bootstrap procedures consistently estimate the marginal sampling
law of \(\widehat{\SR}(m)\) for each fixed \(m\).  If marginal
positive-edge \(p\)-values are valid and the stepdown rejection regions
are monotone, Holm and Romano-Wolf stepdown control familywise error
asymptotically for the fixed menu.
\end{corollary}

\paragraph{Justification.}
This is an assembly of standard results rather than a new theorem.  The
Sharpe functional is Hadamard differentiable on the set
\(m_2-\mu^2\ge\eta\).  A central limit theorem for the sample mean of
\((\bar r_t,\bar r_t^2)\), the delta method, and standard block-bootstrap
consistency for mixing sequences imply marginal Sharpe validity
\citep{kunsch1989jackknife,politis1994stationary,lahiri2003resampling}.
For a fixed finite menu, Holm controls FWER for valid marginal
\(p\)-values, and Romano-Wolf stepdown uses joint max-statistic
resampling under the standard monotonicity requirement
\citep{holm1979simple,romano2005stepwise}.

\begin{remark}[Scope of the validity statement]
The validity statement does not cover long memory, infinite-variance
heavy tails, growing researcher menus, structural breaks, or
nonstationary date returns.  The empirical evaluation reports stationarity
diagnostics and stress simulations, but those outputs are caveats and
diagnostics, not proof of asymptotic validity in those settings.
\end{remark}

\paragraph{Appendix material.}
The worked numerical example, synthetic positive-control table, full
proof details, and secondary robustness artifacts are reported in the
technical appendix.  The main manuscript keeps only the claim boundary
and the evidence needed to evaluate it.

\section{Public benchmark and stress-test evidence}
\label{sec:empirical}

The public experiment is a registry-defined candidate campaign rather than
a single example.  It evaluates the canonical Kenneth French
winner-minus-loser momentum decile benchmark, a broader Kenneth French
decile-rank threshold profile, Kenneth French 25 Size/BM portfolios with
a skipped-momentum signal, and public AQR value, momentum, quality, beta,
and HML-style factor candidates.  The AQR series are economically meaningful
factor-series benchmarks, but they are pre-aggregated series rather than
row-level constituent panels; when the same-date permutation gate cannot
be applied, the candidate is not promoted to a passed panel evaluation.
Unavailable public data sources are reported instead of silently removed.

The generated empirical tables in this section are produced from a
full-campaign root by \path{code/render_manuscript_artifacts.py}.  Table
\ref{tab:candidate-campaign} reports the candidate-level evaluation outcome.
Table \ref{tab:momentum-benchmark} verifies that the canonical WML
candidate matches the direct French winner-minus-loser decile benchmark.
Table \ref{tab:phantom-audit} gives the dependence audit
summary using row \(p\)-values, UVIF, EIR, audit \(p\)-values, and net
Sharpe status for panel candidates.  Table \ref{tab:horizon-effect}
reports the targeted \(L=1,5,21\) horizon-effect audit for the dynamic
Size/BM momentum panel.  The technical appendix reports the full
secondary table set: sampling-boundary diagnostics, same-date
permutation details, HAC bandwidth/prewhitening/fixed-\(b\) robustness,
gate sensitivity, holdout splits, and cost sweeps.

The factor-alpha, data-snooping, stationarity, holdout, and cost tables
are secondary diagnostics, not replacements for the primary date-return
decision.  A raw positive Sharpe that does not survive FF3 plus momentum
alpha, bootstrap robustness, Romano-Wolf/White checks, same-date
placebos, or plausible costs should not be interpreted as evidence of a
tradable edge.  This is the empirical "so what" of the protocol: familiar
public candidates can look attractive at one boundary and fail at the
full decision boundary.

% Generated by render_manuscript_artifacts.py from campaign root code/output_overhaul_20260601
% Main-manuscript subset generated from generated_empirical_artifacts.tex.

The composite campaign table shows that only the canonical French WML
panel clears all 5 percent gates; other positive factor-series evidence
is not treated as a passed panel claim when the row-level placebo is not
applicable.

\begin{table}[t]
\centering
\small
\caption{Registry-defined public candidate evaluation campaign.}
\label{tab:candidate-campaign}
\resizebox{\textwidth}{!}{%
\begin{tabular}{lllllllll}
\toprule
Candidate & status & SR & HAC p+ & Boot p+ & RW p & Perm p & Net SR & Pass 5\% \\
\midrule
french\_momentum\_deciles\_daily\_rank & Available & 0.447 & 7.54e-05 & 2.00e-04 & 2.00e-04 & 1.000 & 0.447 & no \\
french\_momentum\_deciles\_daily\_wml & Available & 0.497 & 5.53e-06 & 2.00e-04 & 2.00e-04 & 1.00e-04 & 0.497 & yes \\
french\_size\_bm\_daily\_dynamic\_momentum & Available & 0.036 & 0.445 & 0.458 & 0.516 & 0.947 & -0.072 & no \\
stooq\_fixed\_etf\_daily\_dynamic\_momentum & Data unavailable &  &  &  &  &  &  & no \\
aqr\_bab\_equity\_factors\_daily & Available & 0.741 & 3.08e-13 & 2.00e-04 & 2.00e-04 &  & 0.741 & no \\
aqr\_qmj\_factors\_daily & Available & 0.536 & 2.26e-05 & 2.00e-04 & 2.00e-04 &  & 0.536 & no \\
aqr\_hml\_devil\_factors\_daily & Available & 0.336 & 0.004 & 0.002 & 0.002 &  & 0.336 & no \\
aqr\_value\_momentum\_everywhere\_monthly & Available & 2.623 & 1.62e-04 & 2.00e-04 & 2.00e-04 &  & 2.623 & no \\
\bottomrule
\end{tabular}%
}
\end{table}

The benchmark validation table shows that the panel implementation
matches the direct French winner-minus-loser construction.

\begin{table}[t]
\centering
\small
\caption{Canonical French momentum benchmark validation.}
\label{tab:momentum-benchmark}
\resizebox{\textwidth}{!}{%
\begin{tabular}{llllllll}
\toprule
Candidate & thr & n & start & end & panel SR & direct WML SR & scale \\
\midrule
french\_momentum\_deciles\_daily\_wml & 0.900 & 26131 & 1926-11-04 & 2026-04-30 & 0.497 & 0.497 & 0.500 \\
\bottomrule
\end{tabular}%
}
\end{table}

The dependence audit table summarizes the boundary directly: row-level
evidence can disappear once same-date redundancy, EIR, placebo evidence,
and turnover-scaled net Sharpe are evaluated together.

\begin{table}[t]
\centering
\small
\caption{Dependence audit summary for panel candidates.}
\label{tab:phantom-audit}
\resizebox{\textwidth}{!}{%
\begin{tabular}{lllllll}
\toprule
Candidate & Row p+ & UVIF & EIR & Audit p+ & Net SR & Status \\
\midrule
french\_momentum\_deciles\_daily\_rank & 1.70e-04 & 1.000 & 1.000 & 1.000 & 0.447 & Fails permutation gate \\
french\_momentum\_deciles\_daily\_wml & 4.85e-04 & 1.000 & 1.000 & 2.00e-04 & 0.497 & Passes composite gate \\
french\_size\_bm\_daily\_dynamic\_momentum & 0.253 & 9.002 & 0.111 & 0.947 & -0.072 & Fails cost gate \\
\bottomrule
\end{tabular}%
}
\par\footnotesize\emph{Note:} EIR is a diagnostic approximation and is not used as a rejection rule.
\end{table}

The horizon table shows that changing the lookback reduces turnover in
the Size/BM stress panel but does not convert it into a robust net-Sharpe
claim.

\begin{table}[t]
\centering
\small
\caption{Horizon-effect audit for the dynamic Size/BM momentum panel.}
\label{tab:horizon-effect}
\resizebox{\textwidth}{!}{%
\begin{tabular}{llllllllll}
\toprule
L & thr & dates & avg names & gross SR & net SR 5bps & turnover & UVIF SR & serial UVIF & HAC p+ \\
\midrule
1 & 0.500 & 3497 & 12.913 & 0.073 & -0.394 & 0.733 & 11.498 & 0.891 & 0.386 \\
5 & 0.500 & 3493 & 13.000 & -0.312 & -0.528 & 0.339 & 12.880 & 0.990 & 0.878 \\
21 & 0.500 & 3477 & 13.000 & 0.058 & -0.050 & 0.169 & 12.121 & 0.933 & 0.412 \\
\bottomrule
\end{tabular}%
}
\end{table}


The campaign separates benchmark validation from broader search
claims.  The canonical French WML candidate is expected to match the
direct Kenneth French winner-minus-loser decile construction and is
reported as a public benchmark check.  This resolves the apparent
\(\widehat{\SR}=0.082\) benchmark problem in earlier artifacts: that
number came from an earlier rank-threshold stress candidate, not from the
canonical WML construction.  In the rendered campaign, the WML panel
Sharpe matches the direct WML Sharpe at approximately 0.50, and the
scale difference is exactly one-half because the panel portfolio holds
equal long and short legs while the direct WML series is winner minus
loser.  The broader rank-threshold profile and the French Size/BM
dynamic-momentum panel are evaluated as stress cases rather than as
canonical anomaly replications.  The AQR
factor-series candidates have positive time-series evidence in the
reported HAC and bootstrap diagnostics, but because they are
pre-aggregated factor returns rather than row-level signal panels, they
cannot be promoted to passed panel-evaluation candidates.

The horizon-effect audit shows that the dynamic Size/BM momentum stress
case fails because of weak signal and execution drag, not because the
canonical momentum benchmark is missing.  At threshold 0.5, \(L=1\)
has very high turnover and a negative 5 bps net Sharpe; \(L=21\) lowers
turnover substantially but still leaves net Sharpe slightly negative.
The empirical Sharpe UVIF relative to row-naive inference is about
11--13 across the three horizons, while the serial-only UVIF is close to
one in this sample.  Thus the dominant inflation channel in this public
stress panel is cross-sectional row multiplication; a longer lookback
does not mechanically make serial dependence the dominant force.
For the dynamic Size/BM panel in Table \ref{tab:phantom-audit}, the
reported EIR is \(0.111\).  Under the equicorrelation diagnostic, this
means \(1-0.111=0.889\), so approximately 89 percent of the row-level
information implied by naive row pooling is redundant.
The dependence audit summary is intentionally artifact-backed rather than
illustrative.  It does not assert that every named factor is a panel
candidate.  Pre-aggregated AQR factor series are excluded from that
table because they do not provide row-level signal-return pairs from
which row \(p\)-values, same-date permutation tests, or panel EIR can be
computed.

\section{Monte Carlo evidence}
\label{sec:simulation}

The Monte Carlo study is designed as a size-and-power audit of the
sampling boundary, not as a search for a favorable strategy.  Each
replication generates a symbol-date return panel, applies the same
date-aggregation rule used in the empirical protocol, and then asks
whether each inference method rejects a positive-edge null at the
nominal 5 percent level.  The base panel is
\[
  r_{t,i}=\mu_t+\beta_i'f_t+\epsilon_{t,i},
  \qquad
  f_t=\Phi f_{t-1}+u_t .
\]
The grid varies serial dependence, cross-sectional dependence, GARCH
volatility, multifactor exposure, and regime switching.  The
high-persistence AR(1) case is reported as such; it is not labeled as a
long-memory design.  The high-volatility GARCH case is not described as
an infinite-variance heavy-tail design.
The null designs set the true edge to zero while varying
cross-sectional correlation and serial memory; the power designs then
increase the true annualized Sharpe on the same dependence-aware
decision boundary.  The important object is the gap between the
row-naive rejection rate and the rejection rates of methods that operate
on date-level returns.  In the displayed null grid, the row-naive test
overrejects around 35--37 percent under dependent nulls, while HAC-delta,
moving-block bootstrap, stationary bootstrap, and Romano-Wolf remain much
closer to nominal size.  In the power grid, row-naive rejection reaches
about 50 percent at a true annualized Sharpe near 0.19, before the
dependence-aware methods show comparable power; that gap is false
precision, not usable evidence.

The size and power experiment compares row-level naive testing,
date-IID bootstrap, HAC-delta inference, moving-block bootstrap,
stationary bootstrap, and Romano-Wolf joint resampling.  The coverage
experiment focuses on IID, moving-block, and stationary bootstrap
intervals across the broader DGP grid.  The purpose is not to rank
methods by power alone.  A method that appears powerful because it
overrejects under the null is not acceptable for confirmatory evaluation.

Table \ref{tab:size} gives the main size evidence.  The row-naive column is the familiar
symbol-date analysis; the HAC-delta, moving-block, stationary, and
Romano-Wolf columns make decisions from the date-return series.  When
row-naive rejection is high in null or near-null designs but the
dependence-aware methods remain near nominal size, the apparent extra
power is a false precision gain from treating same-date rows as
independent evidence.  The technical appendix reports the full DGP grid,
coverage table, power table, and power-curve figure.

% Main-manuscript subset generated from generated_simulation_artifacts.tex.
% Full DGP, coverage, power tables, and power-curve figure are retained in
% sddm_technical_appendix.tex.
The null-size table is the core simulation result: under dependent nulls,
row-naive testing rejects about one third of the time at a nominal
5 percent level, while date-level HAC, block-bootstrap, stationary, and
Romano-Wolf procedures stay much closer to size.

\begin{table}[t]
\centering
\small
\caption{Null rejection rates at nominal 5 percent size.}
\label{tab:size}
\resizebox{\textwidth}{!}{%
\begin{tabular}{lllllll}
\toprule
dgp & date\_iid & hac\_delta & moving\_block & romano\_wolf & row\_naive & stationary \\
\midrule
null\_dep & 0.094 & 0.066 & 0.070 & 0.068 & 0.352 & 0.067 \\
null\_garch & 0.098 & 0.063 & 0.070 & 0.067 & 0.351 & 0.068 \\
null\_iid & 0.060 & 0.060 & 0.061 & 0.060 & 0.062 & 0.062 \\
\bottomrule
\end{tabular}%
}
\end{table}


\section{Discussion and limitations}

The results should be read as a claim-boundary exercise rather than a
strategy discovery exercise.  The framework addresses two common sources
of false precision in trading-panel evidence: same-date cross-sectional
row multiplication and serial dependence in the date-level return
series.  It does not remove the need for economic interpretation,
out-of-sample validation, liquidity analysis, or venue-specific
execution modeling.

The formal validity statement assumes a fixed and finite researcher menu,
stationary date-return series, adequate moments for the Sharpe
delta-method expansion, and block-bootstrap regularity.  The framework
does not cover long-memory processes, infinite-variance returns, growing
researcher menus, structural breaks, or nonstationary date returns
without additional assumptions.  The empirical stationarity, holdout,
factor-alpha, placebo, and cost outputs are therefore diagnostics and
boundary checks, not proofs that every maintained assumption holds in
every market setting.

\section{Data and code availability}

The reproducibility package contains the LaTeX source, generated public
aggregate artifacts, public-data loaders, simulation code, unit tests,
release scripts, and SHA256 provenance.  The public repository is:
{\scriptsize\noindent
\url{https://github.com/DsChauhan08/dependence-aware-evaluation-of-panel-trading-strategy-evidence}\par}
Code is released under Apache-2.0 and manuscript/public aggregate
artifacts under CC-BY 4.0 unless a target venue requires different terms.
Raw third-party data are not
redistributed where source terms do not permit redistribution; the
repository instead provides documented loaders, source links, access
metadata, generated aggregate artifacts, and checksums needed to
reproduce the analysis.  Proprietary predictions, raw trades, private
tickers, production feature definitions, and unrelated workspace
artifacts are excluded.  The technical appendix contains the full
simulation grid, robustness tables, positive-control artifact, registry
notes, and proof details omitted from the main manuscript for length.

\section{Conclusion}

This paper provides a dependence-aware evaluation protocol for panel
trading strategy evidence.  It moves claims from selected symbol-date
rows to date-level portfolio returns, estimates Sharpe uncertainty with
HAC and dependent-bootstrap methods, corrects the fixed researcher menu,
checks same-date placebo and factor-alpha boundaries, and reports
turnover-scaled cost sensitivity.  The Sharpe UVIF converts the familiar
design-effect problem into a variance ratio for the nonlinear Sharpe
functional, and EIR reports how much nominal row-level information
remains nonredundant after same-date dependence is recognized.  For
researchers and allocators, the practical implication is direct:
statistical significance earned by multiplying dependent rows is not
evidence of a tradable edge.

The public campaign shows how this claim boundary changes interpretation.
The canonical French WML candidate validates the pipeline against the
direct winner-minus-loser construction.  By contrast, broader
rank-threshold and dynamic Size/BM panel claims fail dependence-aware,
placebo, researcher-menu, or cost boundaries; in the Size/BM stress
panel, an EIR of \(0.111\) implies that about 89 percent of nominal
row-level information is redundant under the diagnostic approximation.
Pre-aggregated AQR factor series remain useful time-series benchmarks,
but without row-level constituent signals they cannot clear the panel
permutation gate.

The protocol is therefore a reproducible claim boundary rather than a
strategy generator.  It helps separate benchmark validation from
exploratory search and makes the cost of dependence explicit before a
result is promoted as alpha.  Future work should extend the framework to
nonlinear market-impact and capacity models, time-varying
cross-sectional dependence, and richer blocked-permutation designs.

\section*{Author contributions}

Dhananjay S. Chauhan is the sole author of this manuscript and was
responsible for conceptualization, methodology, software, validation,
formal analysis, data curation, writing, and revision.

\section*{Acknowledgments}

This research did not receive any specific grant from funding agencies in
the public, commercial, or not-for-profit sectors.

\section*{Conflict of interest}

The author declares no conflict of interest.

\section*{Ethics statement}

This study uses publicly available aggregate financial data and
simulation data. It does not involve human participants, personal data,
clinical data, animal subjects, or cell lines. Ethics approval was
therefore not required.

\section*{Data availability}

The data used in this study are derived from publicly available financial
datasets, including Kenneth French data and AQR factor data. Cleaned
public aggregate artifacts, source metadata, checksums, and reproduction
instructions are available at:
{\scriptsize\noindent
\url{https://github.com/DsChauhan08/dependence-aware-evaluation-of-panel-trading-strategy-evidence}\par}
Raw third-party data are not redistributed where source terms do not
permit redistribution.

\section*{Code availability}

The code used to generate the empirical results, Monte Carlo simulations,
tables, figures, and manuscript artifacts is available at:
{\scriptsize\noindent
\url{https://github.com/DsChauhan08/dependence-aware-evaluation-of-panel-trading-strategy-evidence}\par}

\bibliographystyle{apalike}
\begin{thebibliography}{99}

\bibitem[Almgren and Chriss, 2001]{almgren2001optimal}
Almgren, R. and Chriss, N. (2001).
\newblock Optimal execution of portfolio transactions.
\newblock {\em Journal of Risk}, 3:5--39.

\bibitem[Andrews, 1991]{andrews1991heteroskedasticity}
Andrews, D. W. K. (1991).
\newblock Heteroskedasticity and autocorrelation consistent covariance matrix estimation.
\newblock {\em Econometrica}, 59(3):817--858.

\bibitem[Andrews and Monahan, 1992]{andrews1992improved}
Andrews, D. W. K. and Monahan, J. C. (1992).
\newblock An improved heteroskedasticity and autocorrelation consistent covariance matrix estimator.
\newblock {\em Econometrica}, 60(4):953--966.

\bibitem[Bailey and Lopez de Prado, 2014]{bailey2014deflated}
Bailey, D. H. and Lopez de Prado, M. (2014).
\newblock The deflated Sharpe ratio: Correcting for selection bias, backtest overfitting, and non-normality.
\newblock {\em Journal of Portfolio Management}, 40(5):94--107.

\bibitem[Bailey et~al., 2014]{bailey2014probability}
Bailey, D. H., Borwein, J. M., Lopez de Prado, M., and Zhu, Q. J. (2014).
\newblock The probability of backtest overfitting.
\newblock {\em Journal of Computational Finance}, 20(4):39--70.

\bibitem[Barras et~al., 2010]{barras2010false}
Barras, L., Scaillet, O., and Wermers, R. (2010).
\newblock False discoveries in mutual fund performance: Measuring luck in estimated alphas.
\newblock {\em Journal of Finance}, 65(1):179--216.

\bibitem[Benjamini and Hochberg, 1995]{benjamini1995controlling}
Benjamini, Y. and Hochberg, Y. (1995).
\newblock Controlling the false discovery rate.
\newblock {\em Journal of the Royal Statistical Society: Series B}, 57(1):289--300.

\bibitem[Benjamini and Yekutieli, 2001]{benjamini2001control}
Benjamini, Y. and Yekutieli, D. (2001).
\newblock The control of the false discovery rate in multiple testing under dependency.
\newblock {\em Annals of Statistics}, 29(4):1165--1188.

\bibitem[Bouchaud et~al., 2008]{bouchaud2008markets}
Bouchaud, J.-P., Farmer, J. D., and Lillo, F. (2008).
\newblock How markets slowly digest changes in supply and demand.
\newblock In {\em Handbook of Financial Markets: Dynamics and Evolution}, pages 57--160.

\bibitem[Carhart, 1997]{carhart1997persistence}
Carhart, M. M. (1997).
\newblock On persistence in mutual fund performance.
\newblock {\em Journal of Finance}, 52(1):57--82.

\bibitem[Cover and Thomas, 2006]{cover2006elements}
Cover, T. M. and Thomas, J. A. (2006).
\newblock {\em Elements of Information Theory}.
\newblock Wiley, 2nd edition.

\bibitem[Fama and French, 1993]{fama1993common}
Fama, E. F. and French, K. R. (1993).
\newblock Common risk factors in the returns on stocks and bonds.
\newblock {\em Journal of Financial Economics}, 33(1):3--56.

\bibitem[Fama and French, 2010]{fama2010luck}
Fama, E. F. and French, K. R. (2010).
\newblock Luck versus skill in the cross-section of mutual fund returns.
\newblock {\em Journal of Finance}, 65(5):1915--1947.

\bibitem[Frazzini et~al., 2015]{frazzini2015trading}
Frazzini, A., Israel, R., and Moskowitz, T. J. (2015).
\newblock Trading costs of asset pricing anomalies.
\newblock Working paper.

\bibitem[Harvey and Liu, 2020a]{harvey2020false}
Harvey, C. R. and Liu, Y. (2020).
\newblock False (and missed) discoveries in financial economics.
\newblock {\em Journal of Finance}, 75(5):2503--2553.

\bibitem[Harvey and Liu, 2020b]{harvey2020census}
Harvey, C. R. and Liu, Y. (2020).
\newblock A census of the factor zoo.
\newblock SSRN working paper, revised October 2020.

\bibitem[Harvey et~al., 2016]{harvey2016and}
Harvey, C. R., Liu, Y., and Zhu, H. (2016).
\newblock ... and the cross-section of expected returns.
\newblock {\em Review of Financial Studies}, 29(1):5--68.

\bibitem[Hasbrouck, 2009]{hasbrouck2009trading}
Hasbrouck, J. (2009).
\newblock Trading costs and returns for U.S. equities.
\newblock {\em Journal of Finance}, 64(3):1445--1477.

\bibitem[Holm, 1979]{holm1979simple}
Holm, S. (1979).
\newblock A simple sequentially rejective multiple test procedure.
\newblock {\em Scandinavian Journal of Statistics}, 6(2):65--70.

\bibitem[Kish, 1965]{kish1965survey}
Kish, L. (1965).
\newblock {\em Survey Sampling}.
\newblock Wiley.

\bibitem[Kiefer and Vogelsang, 2005]{kiefer2005new}
Kiefer, N. M. and Vogelsang, T. J. (2005).
\newblock A new asymptotic theory for heteroskedasticity-autocorrelation robust tests.
\newblock {\em Journal of Econometrics}, 129(1--2):131--162.

\bibitem[Kunsch, 1989]{kunsch1989jackknife}
Kunsch, H. R. (1989).
\newblock The jackknife and the bootstrap for general stationary observations.
\newblock {\em Annals of Statistics}, 17(3):1217--1241.

\bibitem[Lahiri, 2003]{lahiri2003resampling}
Lahiri, S. N. (2003).
\newblock {\em Resampling Methods for Dependent Data}.
\newblock Springer.

\bibitem[Ledoit and Wolf, 2008]{ledoit2008robust}
Ledoit, O. and Wolf, M. (2008).
\newblock Robust performance hypothesis testing with the Sharpe ratio.
\newblock {\em Journal of Empirical Finance}, 15(5):850--859.

\bibitem[Lo, 2002]{lo2002statistics}
Lo, A. W. (2002).
\newblock The statistics of Sharpe ratios.
\newblock {\em Financial Analysts Journal}, 58(4):36--52.

\bibitem[McLean and Pontiff, 2016]{mclean2016does}
McLean, R. D. and Pontiff, J. (2016).
\newblock Does academic research destroy stock return predictability?
\newblock {\em Journal of Finance}, 71(1):5--32.

\bibitem[Moulton, 1986]{moulton1986random}
Moulton, B. R. (1986).
\newblock Random group effects and the precision of regression estimates.
\newblock {\em Journal of Econometrics}, 32(3):385--397.

\bibitem[Novy-Marx, 2014]{novymarx2014predicting}
Novy-Marx, R. (2014).
\newblock Predicting anomaly performance with politics, the weather, global warming, sunspots, and the stars.
\newblock {\em Journal of Financial Economics}, 112(2):137--146.

\bibitem[Opdyke, 2007]{opdyke2007comparing}
Opdyke, J. D. (2007).
\newblock Comparing Sharpe ratios: So where are the p-values?
\newblock {\em Journal of Asset Management}, 8(5):308--336.

\bibitem[Politis and Romano, 1994]{politis1994stationary}
Politis, D. N. and Romano, J. P. (1994).
\newblock The stationary bootstrap.
\newblock {\em Journal of the American Statistical Association}, 89(428):1303--1313.

\bibitem[Politis and White, 2004]{politis2004automatic}
Politis, D. N. and White, H. (2004).
\newblock Automatic block-length selection for the dependent bootstrap.
\newblock {\em Econometric Reviews}, 23(1):53--70.

\bibitem[Petersen, 2009]{petersen2009estimating}
Petersen, M. A. (2009).
\newblock Estimating standard errors in finance panel data sets: Comparing approaches.
\newblock {\em Review of Financial Studies}, 22(1):435--480.

\bibitem[Romano and Wolf, 2005]{romano2005stepwise}
Romano, J. P. and Wolf, M. (2005).
\newblock Stepwise multiple testing as formalized data snooping.
\newblock {\em Econometrica}, 73(4):1237--1282.

\bibitem[Storey, 2002]{storey2002direct}
Storey, J. D. (2002).
\newblock A direct approach to false discovery rates.
\newblock {\em Journal of the Royal Statistical Society: Series B}, 64(3):479--498.

\bibitem[White, 2000]{white2000reality}
White, H. (2000).
\newblock A reality check for data snooping.
\newblock {\em Econometrica}, 68(5):1097--1126.

\end{thebibliography}

\end{document}
